\documentclass[a4paper,11pt]{article}

% ============================================================
%  Le Collé — Algèbre linéaire (compléments de deuxième année)
%  © 2026 Ahmed Amine Ouqlani
% ============================================================

% ---------- Polices ----------
\usepackage[utf8]{inputenc}
\usepackage[T1]{fontenc}
\usepackage[sc]{mathpazo}          % texte et mathématiques : Palatino
\usepackage[scaled=0.88]{helvet}   % titres et repères : linéale
\linespread{1.07}
\usepackage{amsmath,amssymb,amsthm,mathrsfs,stmaryrd}
\usepackage{microtype}

% ---------- Mise en page ----------
\usepackage[a4paper,top=2.3cm,bottom=2.3cm,left=2.4cm,right=2.4cm,
            headheight=14pt,headsep=0.8cm,footskip=1.1cm]{geometry}
\setlength{\parindent}{0pt}
\setlength{\parskip}{0.45em plus 0.15em minus 0.1em}
\raggedbottom
\emergencystretch=1.5em
\clubpenalty=10000 \widowpenalty=10000 \displaywidowpenalty=10000 \brokenpenalty=10000

% ---------- Couleurs : un bleu de nuit, un ocre, des gris ----------
\usepackage{xcolor}
\definecolor{Ink}{HTML}{1C2533}
\definecolor{Bleu}{HTML}{1F3A5F}
\definecolor{BleuMoyen}{HTML}{4A6A8F}
\definecolor{BleuClair}{HTML}{9DB2C9}
\definecolor{Fond}{HTML}{F2F4F7}
\definecolor{Ocre}{HTML}{B07A26}
\definecolor{Brique}{HTML}{A23B2A}
\definecolor{Gris}{HTML}{5E6672}
\definecolor{Filet}{HTML}{C5CCD6}

% ---------- Packages ----------
\usepackage{enumitem}
\setlist{itemsep=0.25em,topsep=0.35em,parsep=0pt}
\setlist[itemize]{label=\textcolor{BleuMoyen}{\small\textbullet},leftmargin=1.4em}
\setlist[enumerate]{leftmargin=1.8em,label=\textsf{\small\color{Bleu}\arabic*.}}
\usepackage{booktabs,array,multicol}
\usepackage{needspace,float}
\usepackage[font=small,labelfont={sf,color=Bleu},labelsep=period,skip=6pt]{caption}
\usepackage{tikz}
\usetikzlibrary{arrows.meta,calc,decorations.markings}
\usepackage{pgfplots}
\pgfplotsset{compat=1.17}
\usepackage{tkz-tab}
\usepackage[most]{tcolorbox}
\usepackage{titlesec,titletoc}
\usepackage{fancyhdr,lastpage,etoolbox}
\usepackage[colorlinks=true,linkcolor=Bleu,urlcolor=Bleu,citecolor=Bleu,
            pdftitle={DS blanc — d après Banque PT — 28 novembre 2026},pdfauthor={Ahmed Amine Ouqlani}]{hyperref}

\numberwithin{equation}{section}
\numberwithin{figure}{section}
\numberwithin{table}{section}

% ---------- Titres ----------
% numéro du chapitre dans le parcours PT (affiché), distinct de la numérotation interne
\newcommand{\grandnum}{\thesection}\newcommand{\ongletcouleur}{Bleu}
% Chapitres (sections) : grand numéro, titre en linéale, filet ocre
\titleformat{\section}[display]
  {\normalfont\sffamily\color{Ink}}
  {{\small\color{Gris}\textls[200]{\MakeUppercase{Chapitre}}}\\[-0.2em]%
   {\fontsize{64}{64}\selectfont\color{BleuClair}\grandnum}}
  {0.4em}
  {\fontsize{25}{30}\selectfont\bfseries}
  [\vspace{0.5em}{\color{Ocre}\rule{3.2cm}{1.4pt}}]
\titlespacing*{\section}{0pt}{0pt}{1.6em}
\titleformat{name=\section,numberless}[display]
  {\normalfont\sffamily\color{Ink}}{}{0pt}{\fontsize{25}{30}\selectfont\bfseries}
  [\vspace{0.5em}{\color{Ocre}\rule{3.2cm}{1.4pt}}]
\titlespacing*{name=\section,numberless}{0pt}{0pt}{1.6em}

\titleformat{\subsection}
  {\needspace{5\baselineskip}\normalfont\sffamily\Large\bfseries\color{Ink}}
  {\color{Bleu}\thesubsection}{0.8em}{}
\titlespacing*{\subsection}{0pt}{1.8em}{0.5em}
\titleformat{\subsubsection}
  {\needspace{5\baselineskip}\normalfont\sffamily\large\bfseries\color{Ink}}
  {\color{Bleu}\thesubsubsection}{0.7em}{}
\titlespacing*{\subsubsection}{0pt}{1.6em}{0.4em}
\titleformat{\paragraph}[runin]{\normalfont\sffamily\bfseries\color{Ink}}{}{0pt}{}[.]
\titlespacing*{\paragraph}{0pt}{0.9em}{0.6em}
\setcounter{secnumdepth}{3}
\setcounter{tocdepth}{2}

% ---------- Table des matières ----------
\titlecontents{section}[0pt]{\addvspace{0.7em}\sffamily\bfseries\color{Ink}}
  {}{}{\hfill\contentspage}[\addvspace{0.2em}]
\titlecontents{subsection}[2.2em]{\small\linespread{1}\selectfont}
  {\makebox[2.6em][l]{\color{Gris}\thecontentslabel}}{}{\titlerule*[0.7pc]{.}\contentspage}

% ---------- En-têtes et pieds de page ----------
\pagestyle{fancy}
\fancyhf{}
\renewcommand{\sectionmark}[1]{\markboth{#1}{}}
\fancyhead[L]{\footnotesize\sffamily\color{Gris}\leftmark}
\fancyhead[R]{\footnotesize\sffamily\scshape\color{Gris}Le Collé}
\fancyfoot[C]{\small\sffamily\color{Ink}\thepage}
\fancyfoot[L]{\scriptsize\sffamily\hyperlink{toc}{\color{Filet}sommaire}}
\renewcommand{\headrulewidth}{0.4pt}
\renewcommand{\headrule}{{\color{Filet}\hrule height\headrulewidth width\headwidth}}
\renewcommand{\footrulewidth}{0pt}
\fancypagestyle{plain}{\fancyhf{}\fancyfoot[C]{\small\sffamily\color{Ink}\thepage}\renewcommand{\headrulewidth}{0pt}}

% ---------- Énoncés ----------
\newtheoremstyle{enonce}{0pt}{0pt}{}{}{\sffamily\bfseries\color{Bleu}}{.}{0.5em}
  {\thmname{#1}\thmnumber{~#2}\thmnote{\ \normalfont\sffamily\color{Ink}(#3)}}
\newtheoremstyle{exemple}{0.9em}{0.9em}{}{}{\sffamily\bfseries\color{BleuMoyen}}{.}{0.5em}
  {\thmname{#1}\thmnumber{~#2}\thmnote{\ \normalfont\sffamily\color{Ink}(#3)}}
\newtheoremstyle{remarque}{0.7em}{0.7em}{}{}{\sffamily\itshape}{.}{0.5em}{\thmname{#1}\thmnote{\ (#3)}}
\theoremstyle{enonce}
\newtheorem{definition}{Définition}[section]
\newtheorem{theoreme}[definition]{Théorème}
\newtheorem{proposition}[definition]{Proposition}
\newtheorem{corollaire}[definition]{Corollaire}
\theoremstyle{exemple}
\newtheorem{exemple}[definition]{Exemple}
\theoremstyle{remarque}
\newtheorem*{remarque}{Remarque}

\tcbset{
  defi/.style={enhanced,breakable,colback=white,colframe=white,boxrule=0pt,arc=0pt,
     borderline west={1.2pt}{0pt}{Bleu},left=11pt,right=4pt,top=3pt,bottom=3pt,
     before skip=1.1em,after skip=1.1em},
  thm/.style={enhanced,breakable,colback=Fond,colframe=Fond,boxrule=0pt,arc=0pt,
     borderline west={2.4pt}{0pt}{Bleu},left=11pt,right=8pt,top=6pt,bottom=6pt,
     before skip=1.2em,after skip=1.2em}}
\tcolorboxenvironment{definition}{defi}
\tcolorboxenvironment{theoreme}{thm}
\tcolorboxenvironment{proposition}{thm}
\tcolorboxenvironment{corollaire}{thm}

\renewcommand{\proofname}{Démonstration}
\renewcommand{\qedsymbol}{\textcolor{BleuClair}{$\blacksquare$}}
\makeatletter
\renewenvironment{proof}[1][\proofname]{\par\pushQED{\qed}\normalfont
  \topsep0.5em\relax\trivlist\item[\hskip\labelsep\sffamily\itshape\color{Gris}#1.]\ignorespaces}
  {\popQED\endtrivlist\@endpefalse}
\makeatother

% ---------- Méthode : filets haut et bas, titre en bandeau ----------
\makeatletter\newcommand{\methlabelfix}{\protected@edef\@currentlabel{\thetcbcounter}}\makeatother
\newtcolorbox[auto counter,number within=section]{methode}[2][]{%
  enhanced,unbreakable,colback=white,colframe=white,boxrule=0pt,arc=0pt,
  borderline north={0.9pt}{0pt}{Bleu},borderline south={0.9pt}{0pt}{Bleu},
  left=4pt,right=4pt,top=8pt,bottom=8pt,before skip=1.4em,after skip=1.4em,
  before upper={\methlabelfix\phantomsection\addcontentsline{met}{metentry}{\protect\makebox[2.6em][l]{\thetcbcounter}#2}%
     {\sffamily\bfseries\small\color{white}\colorbox{Bleu}{\strut\ MÉTHODE~\thetcbcounter\ }}%
     \ \ {\sffamily\bfseries\color{Bleu}#2}\par\smallskip},#1}

% ---------- Repères pédagogiques : un filet de couleur, un mot-clé ----------
\newcommand{\definirbloc}[3]{% nom, couleur, libellé
  \newtcolorbox{#1}[1][]{enhanced,breakable,blanker,
    borderline west={1.4pt}{0pt}{#2},left=10pt,right=2pt,top=1pt,bottom=1pt,
    before skip=1em,after skip=1em,
    before upper={{\sffamily\bfseries\footnotesize\color{#2}\textls[120]{\MakeUppercase{#3}}}\quad},##1}}
\definirbloc{reflexe}{Bleu}{Réflexe}
\definirbloc{astuce}{Ocre}{Astuce}
\definirbloc{attention}{Brique}{Attention}
\definirbloc{enpratique}{BleuMoyen}{En pratique}
\definirbloc{plusloin}{Gris}{Pour aller plus loin}
\definirbloc{strategie}{Bleu}{Stratégie}
\definirbloc{lecture}{BleuMoyen}{Lire le schéma}
% compatibilité avec les parties reprises
\newenvironment{garde}{\begin{attention}}{\end{attention}}
\newenvironment{complementHP}{\begin{plusloin}}{\end{plusloin}}
\newenvironment{interpretation}{\begin{enpratique}}{\end{enpratique}}

% ---------- Exercices, corrigés ----------
\newcounter{exo}[section]
\renewcommand{\theexo}{\thesection.\arabic{exo}}
\newcommand{\niveau}[1]{{\color{Ocre}\footnotesize
  \ifcase#1\or$\bullet\circ\circ\circ\circ$\or$\bullet\bullet\circ\circ\circ$\or$\bullet\bullet\bullet\circ\circ$%
  \or$\bullet\bullet\bullet\bullet\circ$\or$\bullet\bullet\bullet\bullet\bullet$\fi}}
\newcommand{\serie}[1]{\par\needspace{8\baselineskip}\vspace{1.5em}%
  \noindent{\sffamily\bfseries\color{Bleu}#1}\par\nobreak\vspace{-0.4em}%
  \noindent{\color{Filet}\rule{\linewidth}{0.4pt}}\par\nobreak\vspace{0.1em}}
\newenvironment{exo}[3]{%
  \par\needspace{5\baselineskip}\vspace{0.75em}\refstepcounter{exo}\label{e:#1}%
  \def\exolab{#1}%
  \noindent{\sffamily\bfseries\color{Bleu}Exercice~\theexo}\quad\niveau{#2}%
  \hfill{\footnotesize\sffamily\color{Gris}#3}\par\nobreak\vspace{0.15em}}%
  {\par\nopagebreak\noindent\hfill{\scriptsize\sffamily\hyperref[c:\exolab]{\color{Gris}corrigé p.~\pageref{c:\exolab}}}\par}
\newenvironment{corrige}[1]{%
  \par\needspace{6\baselineskip}\vspace{1em}\phantomsection\label{c:#1}%
  \noindent{\sffamily\bfseries\color{Bleu}Exercice~\ref{e:#1}}\quad
  {\scriptsize\sffamily\hyperref[e:#1]{\color{Gris}énoncé p.~\pageref{e:#1}}}\par\nobreak\vspace{0.2em}}%
  {\hfill\textcolor{BleuClair}{$\square$}\par}
\newcommand{\etape}[1]{\par\smallskip\noindent{\sffamily\bfseries\small\color{BleuMoyen}#1.}\ }
\newenvironment{solution}{\par\noindent{\sffamily\bfseries\small\color{BleuMoyen}Solution.}\ }{\hfill\textcolor{BleuClair}{$\square$}\par}
\newcommand{\chercher}{\par\smallskip\noindent{\color{Filet}\rule[0.5ex]{1.2cm}{0.4pt}}\
  {\footnotesize\sffamily\itshape\color{Gris}laisser chercher l'étudiant avant de poursuivre}\par\smallskip}

% ---------- Guide de séance ----------

\newcommand{\marge}[1]{\leavevmode\marginpar{\raggedleft\scriptsize\sffamily\color{Gris}\textls[100]{\MakeUppercase{#1}}}\ignorespaces}
\newcommand{\Etape}[4]{%
  \subsection[Étape #1 --- #3 {\color{Gris}(#2)}]{#3}\label{etape:#1}%
  \noindent{\small\sffamily\color{Gris}{\color{Ocre}\bfseries Étape #1}\quad·\quad #2\quad·\quad #4}\par\medskip}
\newcommand{\question}[1]{\par\smallskip\noindent{\sffamily\bfseries\small\color{BleuMoyen}Question.}\ \textit{#1}\par\smallskip}
\newcommand{\rappel}[1]{\par\smallskip\noindent\phantomsection\addcontentsline{rap}{rapentry}{#1}%
  {\sffamily\bfseries\scriptsize\colorbox{Fond}{\color{Gris}\strut RAPPEL}}\ {\sffamily\small\color{Gris}#1.}\ }
\newcommand{\pourquoi}{\par\noindent{\sffamily\bfseries\small\color{Ocre}Pourquoi ?}\ }

% Ouverture de chapitre : \ouverture{objectifs}{prérequis}{durée}
\newcommand{\ouverture}[3]{%
  \begin{tcolorbox}[enhanced,colback=Fond,colframe=Fond,arc=0pt,boxrule=0pt,left=12pt,right=12pt,top=10pt,bottom=10pt,
     before skip=1.2em,after skip=1.6em]
  \begin{minipage}[t]{0.56\linewidth}
  {\sffamily\bfseries\small\color{Bleu}OBJECTIFS}\par\vspace{0.3em}\small #1
  \end{minipage}\hfill
  \begin{minipage}[t]{0.38\linewidth}
  {\sffamily\bfseries\small\color{Bleu}PRÉREQUIS}\par\vspace{0.3em}\small #2
  \ifx\relax#3\relax\else\par\vspace{0.6em}{\sffamily\bfseries\small\color{Bleu}DURÉE}\par\vspace{0.2em}\small #3\fi
  \end{minipage}
  \end{tcolorbox}}
% Fin de chapitre
\newenvironment{essentiel}{\par\needspace{12\baselineskip}\vspace{1.5em}%
  \begin{tcolorbox}[enhanced,unbreakable,colback=white,colframe=white,boxrule=0pt,arc=0pt,
   borderline north={1.4pt}{0pt}{Ocre},left=2pt,right=2pt,top=8pt,bottom=4pt]
  {\sffamily\bfseries\color{Ocre}L'ESSENTIEL DU CHAPITRE}\par\medskip}{\end{tcolorbox}}

% ---------- Notations ----------
\newcommand{\K}{\mathbb{K}}
\newcommand{\R}{\mathbb{R}}
\newcommand{\C}{\mathbb{C}}
\newcommand{\N}{\mathbb{N}}
\newcommand{\Z}{\mathbb{Z}}
\newcommand{\Mn}[1]{\mathcal{M}_{n}(#1)}
\newcommand{\GLn}[1]{\mathrm{GL}_n(#1)}
\newcommand{\LE}{\mathcal{L}(E)}
\newcommand{\rg}{\operatorname{rg}}
\newcommand{\vect}{\operatorname{Vect}}
\newcommand{\Mat}{\operatorname{Mat}}
\newcommand{\Id}{\operatorname{Id}}
\newcommand{\sh}{\operatorname{sh}}
\newcommand{\ch}{\operatorname{ch}}
\renewcommand{\th}{\operatorname{th}}
\DeclareMathOperator{\argsh}{argsh}
\newcommand{\vv}[1]{\overrightarrow{#1}}
\newcommand{\ve}[1]{\vec{#1}}
\newcommand{\ient}[2]{\llbracket #1,#2\rrbracket}
\newcommand{\dt}[1]{\left|\begin{matrix}#1\end{matrix}\right|}
\newcommand{\dis}{\displaystyle}

% ---------- Figures ----------
\pgfplotsset{trig format plots=rad}
\pgfplotsset{courbe/.style={axis lines=middle,axis equal image,
  axis line style={Gris,-{Stealth[length=5pt]}},
  xlabel={$x$},ylabel={$y$},ticklabel style={font=\scriptsize\sffamily,color=Gris},
  every axis x label/.style={at={(ticklabel* cs:1)},anchor=west,font=\small},
  every axis y label/.style={at={(ticklabel* cs:1)},anchor=south,font=\small},
  samples=300,smooth,width=0.6\linewidth}}
\tikzset{arcetude/.style={Bleu,line width=1.25pt},
         reste/.style={BleuClair,line width=1pt},
         tgte/.style={Ocre,line width=0.9pt},
         asy/.style={Gris,dashed,line width=0.6pt},
         pt/.style={circle,fill=Ink,inner sep=1.2pt}}

% ---------- Navigation ----------
\makeatletter
\newcommand*\l@rapentry{\@dottedtocline{1}{0em}{0em}}
\newcommand*\l@metentry{\@dottedtocline{1}{0em}{0em}}
\newcommand{\listemethodes}{\@starttoc{met}}
\newcommand{\listerappels}{\@starttoc{rap}}
\makeatother
% onglets de chapitre sur le bord droit
\usepackage{eso-pic}
\newif\ifonglets \ongletsfalse
\newcommand{\dessineronglet}{%
  \edef\ongcol{\ongletcouleur}%
  \edef\ongnum{\grandnum}%
  \pgfmathsetmacro{\ongy}{3.2+2.3*(\ongnum-1)}%
  \begin{tikzpicture}[remember picture,overlay]
  \fill[\ongcol] ([xshift=-0.75cm,yshift=-\ongy cm]current page.north east)
     rectangle ([yshift=-\ongy cm-1.9cm]current page.north east);
  \node[text=white,font=\sffamily\bfseries\large] at ([xshift=-0.37cm,yshift=-\ongy cm-0.95cm]current page.north east)
     {\ongnum};
  \end{tikzpicture}}
\AddToShipoutPictureBG{\ifonglets\ifnum\value{section}>0 \dessineronglet\fi\fi}
% blocs de navigation et d'enseignement
\definirbloc{commencer}{Gris}{Si vous commencez ici}
\definirbloc{horsprogramme}{Brique}{Hors programme}
\newcommand{\HP}{\,{\sffamily\bfseries\scriptsize\textcolor{Brique}{[HP]}}}
\newcommand{\source}[1]{{\footnotesize\sffamily\color{Gris}\,(#1)}}
\newcommand{\coupdepouce}[1]{\par\vspace{-0.2em}\noindent{\footnotesize\sffamily\color{Gris}%
  \textbf{Coup de pouce} (à donner si l'étudiant bloque) : #1}\par}
\newcommand{\pointarret}[1]{\par\vspace{1em}\noindent{\color{Filet}\rule[0.5ex]{\linewidth}{0.4pt}}\par\nobreak
  \noindent{\footnotesize\sffamily\color{Ocre}\bfseries POINT D'ARRÊT POSSIBLE}\quad{\footnotesize\sffamily\color{Gris}#1}\par\vspace{0.6em}}
\newcommand{\demarche}[1]{\noindent{\sffamily\bfseries\small\color{BleuMoyen}Démarche.}\ \textit{#1}\par\smallskip}

% ---------- Corrigés : idées, astuces, pièges, à retenir ----------
\newcommand{\choix}[1]{\par\noindent{\sffamily\bfseries\small\color{BleuMoyen}Choix de la méthode.}\ #1\par\smallskip}
\newcommand{\gain}[1]{\par\noindent{\sffamily\bfseries\small\color{Ocre}Astuce.}\ #1\par}
\newcommand{\piege}[1]{\par\noindent{\sffamily\bfseries\small\color{Brique}Piège.}\ #1\par}
\newtcolorbox{retenirbox}{enhanced,breakable,colback=Fond,colframe=Fond,arc=0pt,boxrule=0pt,
  borderline west={2pt}{0pt}{Ocre},left=8pt,right=6pt,top=4pt,bottom=4pt,before skip=0.7em,after skip=0.5em}
\newcommand{\retenir}[1]{\begin{retenirbox}{\sffamily\bfseries\footnotesize\color{Ocre}À RETENIR}\quad #1\end{retenirbox}}
\newcommand{\TD}[1]{{\color{Ocre}\bfseries TD\,#1}\ ·\ }
\newcommand{\q}[1]{\par\smallskip\noindent{\sffamily\bfseries\color{Bleu}#1}\ }
\providecommand{\iddots}{\mathinner{\mkern1mu\raise1pt\hbox{.}\mkern2mu\raise4pt\hbox{.}\mkern2mu\raise7pt\hbox{.}\mkern1mu}}


\fancyfoot[L]{}
\newcommand{\Partie}[1]{\par\needspace{5\baselineskip}\vspace{1.1em}\noindent
  {\sffamily\bfseries\large\color{Bleu}#1}\par\vspace{2pt}{\color{Filet}\rule{\linewidth}{0.4pt}}\par\smallskip}
\newcommand{\Q}[1]{\par\needspace{3\baselineskip}\medskip\noindent{\sffamily\bfseries\color{Bleu}#1}\enspace}
\setlist[enumerate,1]{label=\arabic*.,leftmargin=1.6em,itemsep=3pt,topsep=3pt}
\newcommand{\bandeau}[2]{%
\begin{tikzpicture}[remember picture,overlay]
\fill[Bleu] (current page.north west) rectangle ([yshift=-4.2cm]current page.north east);
\node[anchor=north west,text=white,font=\sffamily\footnotesize,inner sep=0] at ([xshift=2.2cm,yshift=-1.1cm]current page.north west)
  {\textls[120]{MATHÉMATIQUES · CPGE PT · DS BLANC}};
\node[anchor=north west,text=white,font=\sffamily\bfseries\LARGE,inner sep=0] at ([xshift=2.2cm,yshift=-1.75cm]current page.north west)
  {DS blanc --- d'après la Banque PT};
\node[anchor=north west,text=white,font=\sffamily\normalsize,inner sep=0] at ([xshift=2.2cm,yshift=-2.75cm]current page.north west)
  {#1};
\node[anchor=north east,text=white,font=\sffamily\scshape\large,inner sep=0] at ([xshift=-2.2cm,yshift=-1.05cm]current page.north east)
  {Le Collé};
\node[anchor=north east,text=white,font=\sffamily\small,inner sep=0] at ([xshift=-2.2cm,yshift=-2.8cm]current page.north east)
  {#2};
\end{tikzpicture}
\vspace*{2.9cm}}
\newcommand{\Source}[1]{\par\vspace{-0.3em}\noindent{\footnotesize\sffamily\color{Gris}#1}\par\smallskip}
\usetikzlibrary{patterns}
\fancyhead[L]{\footnotesize\sffamily\color{Gris}DS blanc · d'après la Banque PT · énoncé}
\providecommand{\ch}{\operatorname{ch}}
\begin{document}
\pagestyle{fancy}\thispagestyle{plain}
\bandeau{Samedi 28 novembre 2026 · préparation au DS du lundi 30 novembre}{Énoncé · 4 h}

\begin{tcolorbox}[enhanced,colback=Fond,colframe=Fond,arc=0pt,boxrule=0pt,left=10pt,right=10pt,top=8pt,bottom=8pt]
\small Sujet composé d'extraits de sujets officiels de la Banque PT (filière PT), choisis pour couvrir le programme du DS :
réduction, séries numériques, équations différentielles et séries entières. Les énoncés sont reformulés et légèrement adaptés
(les questions sur les intégrales généralisées ont été retirées) ; la source de chaque exercice est indiquée. Les quatre exercices
sont indépendants. Calculatrice interdite. La qualité de la rédaction et la précision des justifications seront prises en compte.
\end{tcolorbox}

\Partie{Exercice 1 — Réduction de quelques matrices}
\Source{D'après Banque PT 2026, Mathématiques A, partie A.}

On considère les matrices
\[M_1=\begin{pmatrix}2&6\\1&1\end{pmatrix},\qquad M_2=\begin{pmatrix}3&-2\\1&1\end{pmatrix},\qquad M_3=\begin{pmatrix}5&-4\\1&1\end{pmatrix}.\]
\begin{enumerate}
\item Rappeler la caractérisation des matrices trigonalisables à l'aide du polynôme caractéristique.
\item $M_1$ est-elle diagonalisable dans $\mathcal M_2(\C)$ ? dans $\mathcal M_2(\R)$ ? Déterminer, si possible, une matrice inversible $P_1$ et une
matrice diagonale $D_1$ telles que $M_1=P_1D_1P_1^{-1}$ (le calcul de $P_1^{-1}$ n'est pas demandé).
\item Mêmes questions pour $M_2$.
\item (a) $M_3$ est-elle diagonalisable dans $\mathcal M_2(\R)$ ? Est-elle trigonalisable dans $\mathcal M_2(\R)$ ?\\
(b) On note $f_3$ l'endomorphisme de $\R^2$ canoniquement associé à $M_3$. Déterminer une base de $\R^2$ dans laquelle la matrice de
$f_3$ est $T_3=\begin{pmatrix}3&1\\0&3\end{pmatrix}$.
\end{enumerate}
On s'intéresse désormais aux matrices $M(a,b)=\begin{pmatrix}a&b\\1&1\end{pmatrix}$, où $(a,b)\in\R^2$.
\begin{enumerate}\setcounter{enumi}{4}
\item À quelle condition nécessaire et suffisante la matrice $M(a,b)$ est-elle inversible ?
\item Calculer le discriminant $\Delta(a,b)$ du polynôme caractéristique de $M(a,b)$.
\item (a) Justifier que si $\Delta(a,b)=0$, alors $M(a,b)$ n'est pas diagonalisable dans $\mathcal M_2(\C)$.\\
(b) Lorsque $\Delta(a,b)\neq0$, $M(a,b)$ est-elle diagonalisable ? On précisera si c'est dans $\mathcal M_2(\R)$ et/ou dans $\mathcal M_2(\C)$.
\item Représenter dans un repère du plan, avec les coordonnées $(a,b)$ : l'ensemble des points tels que $M(a,b)$ n'est pas inversible ;
celui des points tels que $M(a,b)$ est diagonalisable dans $\mathcal M_2(\R)$ ; celui des points tels que $M(a,b)$ est diagonalisable dans
$\mathcal M_2(\C)$ mais pas dans $\mathcal M_2(\R)$ ; celui des points tels que $M(a,b)$ n'est pas diagonalisable dans $\mathcal M_2(\C)$ ; enfin les trois
points correspondant à $M_1$, $M_2$, $M_3$.
\end{enumerate}

\Partie{Exercice 2 — Une série alternée de logarithmes}
\Source{D'après Banque PT 2026, Mathématiques C, préambule et partie I (questions sur les intégrales retirées).}

Pour tout entier $n\geq2$, on pose $u_n=\ln\left(1+\dfrac{(-1)^n}n\right)$.
\begin{enumerate}
\item Rappeler le développement limité à l'ordre $2$ en $0$ de $h\mapsto\ln(1+h)$.
\item Déterminer la limite de la suite $(u_n)_{n\geq2}$.
\item Soient $\sum a_n$ et $\sum b_n$ deux séries convergentes et $\lambda\in\R$. Montrer que la série de terme général $a_n+\lambda b_n$ converge.
\item Étudier la convergence et la convergence absolue de la série $\sum_{n\geq2}u_n$.
\item Pour tout entier $N\geq2$, on pose $S_N=\sum_{k=2}^Nu_k$.
\begin{enumerate}[label=(\alph*)]
\item Montrer que, pour tout entier $N\geq1$, $S_{2N}=\ln\left(\dfrac{2N+1}{2N}\right)$.
\item Pour tout entier $n\geq1$, comparer $u_{2n+1}$ et $u_{2n}$.
\item En déduire, pour tout entier $N\geq2$, la valeur de $S_{2N-1}$.
\item Calculer $\displaystyle\sum_{n=2}^{+\infty}u_n$.
\end{enumerate}
\end{enumerate}

\Partie{Exercice 3 — Reste de la série exponentielle, irrationalité de $e$}
\Source{D'après Banque PT 2014, Mathématiques C, partie I.}

Pour $t\in\R$ et $n\in\N$, on pose $R_n(t)=e^t-\displaystyle\sum_{k=0}^n\frac{t^k}{k!}$.
\begin{enumerate}
\item
\begin{enumerate}[label=(\alph*)]
\item Montrer que $R_n$ est solution sur $\R$ de l'équation différentielle $(E)$ : $y'-y=\dfrac{t^n}{n!}$.
\item Donner la solution générale de l'équation homogène $(E_0)$ associée à $(E)$.
\item Résoudre $(E)$ en exprimant les solutions à l'aide d'une intégrale, sans chercher à la calculer.
\item En déduire une expression de $R_n(t)$ à l'aide d'une intégrale, pour tout réel $t$ et tout entier $n$.
\item Montrer que, pour tout réel $t\geq0$, $0\leq R_n(t)\leq\dfrac{e^t\,t^{n+1}}{(n+1)!}$, et que pour tout réel $t<0$, $|R_n(t)|\leq\dfrac{|t|^{n+1}}{(n+1)!}$.
\item À l'aide de ce qui précède, retrouver le développement en série entière de la fonction exponentielle en $0$, en précisant son
rayon de convergence.
\end{enumerate}
\item On considère les suites définies pour $n\geq1$ par $u_n=\displaystyle\sum_{k=0}^n\frac1{k!}$ et $v_n=u_n+\dfrac1{n\,n!}$.
\begin{enumerate}[label=(\alph*)]
\item Montrer que $(u_n)$ et $(v_n)$ sont strictement monotones, et donner leur sens de variation.
\item Montrer que $(u_n)$ et $(v_n)$ sont adjacentes, de limite $e$.
\item En déduire que, pour tout entier $q\geq1$, $u_q<e<v_q$.
\item On suppose qu'il existe deux entiers naturels non nuls $p$ et $q$ tels que $e=\frac pq$. En multipliant la double inégalité précédente
par $q!$, montrer que c'est impossible. Conclure.
\item On pose $w_n=\dfrac1n\displaystyle\sum_{k=1}^nu_k$ pour $n\geq1$. Soit $\varepsilon>0$.
\begin{enumerate}[label=(\roman*)]
\item Montrer qu'il existe un rang $n_0$ tel que, pour tout $n\geq n_0$, $|u_n-e|\leq\frac\varepsilon2$.
\item Montrer qu'il existe un rang $n_1$ tel que, pour tout $n\geq n_1$, $|w_n-e|\leq\varepsilon$.
\item Quelle est la limite de la suite $(w_n)$ ?
\end{enumerate}
\end{enumerate}
\item On considère la suite $(e_n)_{n\geq2}$ définie par $e_n=\left(1-\dfrac1n\right)^{-n}$. Quelle est sa limite ?
\end{enumerate}

\Partie{Exercice 4 — Développement en série entière de la tangente}
\Source{D'après Banque PT 2022, Mathématiques C, partie II.}

Soit $R$ un réel strictement positif.
\begin{enumerate}
\item
\begin{enumerate}[label=(\alph*)]
\item Soient $f$ et $g$ deux fonctions développables en série entière sur $]-R,R[$ : pour tout $x\in\,]-R,R[$,
$f(x)=\sum_{n=0}^{+\infty}a_nx^n$ et $g(x)=\sum_{n=0}^{+\infty}b_nx^n$, où les $a_n$ et $b_n$ sont réels. Rappeler la formule du produit de Cauchy
donnant le produit des séries entières $\sum a_nx^n$ et $\sum b_nx^n$. Que peut-on dire du rayon de convergence de ce produit ?
\item Appliquer la question précédente pour exprimer, pour tout $x\in\,]-R,R[$, le développement en série entière de $\big(f(x)\big)^2$.
\end{enumerate}
\item On suppose désormais que $f(0)=0$ et que, pour tout $x\in\,]-R,R[$, $f'(x)=1+\big(f(x)\big)^2$.
\begin{enumerate}[label=(\alph*)]
\item En déduire la relation de récurrence vérifiée, pour tout entier naturel $n$, par les coefficients $a_n$.
\item Démontrer, par récurrence forte, que $a_{2p}=0$ pour tout entier naturel $p$. Que peut-on en déduire pour la parité de $f$ et pour
son développement en série entière ?
\item Montrer que $f'(0)=1$, puis calculer $a_1$, $a_3$, $a_5$ et $a_7$.
\item On admet que la fonction tangente est développable en série entière sur $\left]-\frac\pi2,\frac\pi2\right[$. Donner le développement limité de
$f$ à l'ordre $7$ en $0$, puis en déduire celui de la fonction tangente au même ordre.
\end{enumerate}
\end{enumerate}

\vfill
\begin{center}\sffamily\small\color{Gris}Fin de l'énoncé\end{center}
\end{document}
